\section{Scientific limitations and uncertainty}\label{sec:limits}
Each limitation is tied to the result it affects. Where it was tested, the test is named.

\subsection{Climate measurement}
\begin{itemize}
\item \textbf{Monthly data hide hot days.} Heat exposure is a lower bound. Rice and soybean score near zero on heat
largely for this reason: their thresholds are defined for short exposures that monthly means cannot show. The
direction of bias is to understate heat vulnerability, most for crops with high thresholds.
\item \textbf{The water ratio has no soil storage or runoff, and treats irrigation as unlimited.}
  \begin{itemize}
  \item It overstates stress for crops grown on stored soil moisture. Rabi chickpea in India is grown on residual
  monsoon moisture, and it sets the water reference.
  \item It understates stress where irrigation water is itself at risk.
  \end{itemize}
  The index therefore assumes irrigation keeps flowing. This matters most for rice and for Telangana.
\item \textbf{Station coverage.} In cell-months without stations, CRU-TS returns climatology, which damps rainfall
variability.
\end{itemize}

\subsection{Crop parameters}
The ledger (Appendix~\ref{app:ledger}) labels every value as published, derived or proxy, and
Section~\ref{sec:parevidence} audits \parNRows\ register rows: \parNPublished\ published, \parNDerived\ derived,
\parNBorrowed\ borrowed, \parNAssumption\ assumptions and \parNUnsupported\ unsupported (the rapeseed heat slope; the
rice heat slope is also unsupported as a yield slope). The ``tested'' ranges below are exploratory perturbations, not
literature-supported uncertainty ranges.
\begin{itemize}
\item \textbf{New legumes are fitted.} Their thresholds and slopes are fits to chamber or pot experiments on one
to three genotypes. Chamber exposure is continuous and without recovery days, so these values probably overstate
field sensitivity. They also sit below model-calibrated thresholds such as those used for wheat, maize and soybean.
Exploratory perturbation: $T_{\text{crit}}+3$\,\si{\celsius} and slopes $\times0.7$ (the direction is supported, the
size is not); wheat stays first.
\item \textbf{Common bean.} The threshold comes from one heat-susceptible Andean cultivar. It is applied to FAO
``beans, dry'', which in India and Myanmar include tropical \emph{Vigna} pulses. Because bean in India sets the heat
reference for all crops, this choice scales every heat exposure; with the India and Myanmar bean pairs removed the
heat reference falls by a factor of \scrHeatRatio\ (Section~\ref{sec:parevidence}). Exploratory perturbation: $+2$\,\si{\celsius}.
\item \textbf{Rice.} FAO gives no rice $K_y$ (FAO-33 classes rice only as ``high'' sensitivity, like maize), so
maize's 1.25 is a proxy. Rice's heat slope is a spikelet-fertility slope, not a yield slope, and is unsupported as
one. Exploratory perturbation: $K_y$ 1.0 and 1.5.
\item \textbf{Rapeseed-mustard.} The threshold is converted from a daily-maximum threshold with an assumed diurnal
range. No published heat slope was found, so wheat's is borrowed and unsupported; it is perturbed from 0.07 to 0.18. The FAO-56 Rev.1
$K_y$ of 0.90 is tabulated without a stated source. Parameters from \emph{B.~napus} probably overstate the heat
sensitivity of Indian mustard (\emph{B.~juncea}).
\item \textbf{Inaccessible sources.} Not read in full: \citet{prasad2002bean}, \citet{prasad2003peanut} and
\citet{morrison2002} (closed access). Their values are taken from abstracts and from \citet{boote2005}.
\end{itemize}

\subsection{Maps, calendars and panels}
\begin{itemize}
\item \textbf{Crop maps are from 2010 (SPAM).} Areas have since moved; Telangana's paddy area is about three times
the 2010 map.
\item \textbf{Calendars.} Calendars are partly gap-filled. Chickpea, lentil and rapeseed use the cool-season rule,
whose plain-Sacks alternative is tested.
\item \textbf{Panels.} Panels sample the main producers (half to nine-tenths of world output), not all countries.
Pigeon pea (tur) and moong, major pulses in South India, are not among the nine crops.
\end{itemize}

\subsection{Nutrition}
\begin{itemize}
\item \textbf{Supply, not intake.} National supply is not individual intake, and composition is one value per item.
Whole-grain and primary-weight compositions overstate zinc and iron.
\item \textbf{Requirements.} RNIs compared with mean supply overstate shortfall at the individual level. The iron
result changes with the bioavailability assumption (\nIronShort\ countries short under the main requirement, none
under the alternative).
\item \textbf{Shock accounting.} There is no price or trade response, so the upper bound is a bound, not a forecast.
\end{itemize}

\subsection{Future}
\begin{itemize}
\item \textbf{One model, one scenario.} The future results show direction under MRI-ESM2-0 SSP2-4.5, not the range
of plausible futures. Regional rainfall changes in particular are model-dependent.
\item \textbf{Fixed area, calendars, irrigation and diets}, with no CO$_2$ fertilisation and no adaptation. Elevated
CO$_2$ would raise yields of C$_3$ crops but lower their zinc, iron and protein \citep{myers2014}. The yield effect is
not included. The concentration effect is not in the main results either; it is reported only as a separate
experimental stress test with partial coverage (Section~\ref{sec:co2}) and must not be read as a forecast.
\item \textbf{Response methods disagree.} The literature and empirical future yield responses disagree for about
half the pairs, so the future nutrient numbers are conditional screening estimates.
\end{itemize}

\subsection{Index design}
\begin{itemize}
\item \textbf{Our combination rules.} Pairing, equal nutrient weights, normalisation by the maximum, and equal or
production weighting across countries are our design choices. Robustness tests show their effect but cannot show
that they are right.
\item \textbf{Relative values.} F4 values are meaningful as relative positions within this set of crops and
countries. Adding crops changes the reference and hence every value, as the four-to-nine-crop comparison showed.
\end{itemize}
